\section{Discussion and Managerial Insights}
\label{sec:discussion}

\subsection{Prediction quality, allocation quality and execution performance}

The two study families support different conclusions. The LOS campaign links
learned stay estimates to session planning and shows that session granularity
and GA mutation settings can alter observed rankings. The room campaign
shows that reaching an exact predicted allocation optimum, or improving
predicted overtime on a larger instance, does not guarantee more completed
cases under hidden durations. Execution evaluates a different quantity from
the surrogate search score. The oracle comparison further shows that more
accurate duration information can increase throughput and overtime together
when the objective prioritises case starts. Prediction MAE, predicted cost,
completions and overtime should therefore be reported separately.

Reactive replanning supports every solver, yet produces gains, ties and
losses. It optimises the visible state with fixed sequencing and reacts only
to closure/reopening, rather than optimising future execution scenarios.
This is a plausible explanation for mixed effects, not evidence that
rescheduling is inherently harmful. Alternative triggers and sequencing
rules require separate experiments.

\subsection{Implications for planning and resource management}

Managers need to assess throughput and overtime jointly. On the 15-case
historical day, lower reactive overtime for tabu and the hybrid accompanies
fewer completed cases; on the 25-case day, overtime falls without changing
completion counts. A recommendation based on overtime alone would conflate
these outcomes. Delay also has a selected population: only completed cases
with an initial assignment contribute, so a method can change both the
average and its denominator. Utilisation excludes turnover and retains
closed room-hours in its denominator; it is not interchangeable with an
availability-adjusted utilisation measure.

Session structure is itself a planning choice: dividing unchanged minutes
into smaller blocks changes the assignment landscape. The LOS hybrid's
5.70\% lower-bound gap in the split-session experiment supports further
comparison under a common budget, not a universal recommendation. Point and
upper-endpoint LOS should remain explicit alternatives with cost and
variability trade-offs. Their relative value depends on workload, weights,
noise rules and bed-accounting conventions.

The capacity-aware date-recommendation ambition is only partially represented. Daily room
experiments allocate a fixed workload and do not propose admission dates,
carry a waiting list across days or enforce ward capacity. The aggregate
model represents beds but has no minute-level execution. Combining their
findings does not establish a validated end-to-end admission recommendation
system. Surgeon availability, staff, equipment, clinical priority and
pre-existing ward occupancy need their own data and constraints.

\subsection{Robustness, generalisation and remaining evidence gaps}

The room campaign uses three solver seeds and only three historical days,
plus two generated samples per size. Seeds repeat a workload and deterministic
baseline repetitions add no workload diversity. Generated rows retain
procedure--duration pairing but omit day-level correlations. The 1.1 target
load and single-room closure do not characterise light-load performance or
constant proportional disruption severity. Equal per-request evaluation
budgets differ from equal policy-wide computation; concurrent 300-case timing
and paused simulation time further limit operational latency conclusions.
Larger held-out workload samples, several loads and budgets, and alternative
disruption patterns are needed before claiming general superiority.

For LOS, a constant-width conformal interval does not capture patient-specific
heteroscedasticity; uniform in-interval noise omits outside-tail events and
clipping changes its mean. Coverage, target units and feature availability
require saved-model verification. About 14\% of cases lack matched predictions,
with fallback depending on the selected workload. Inclusive total LOS differs
from the optimiser's surgery-day-plus-stay occupancy convention; horizon
truncation omits downstream demand. The matched-only variant removes cases,
so its improvement cannot isolate matching quality.

Learned room-model evaluation must pin the current dataset, category mapping,
selected artefact and calibration sample, then provide held-out predictions
for the exact room-occupancy target. Existing median/oracle room experiments
and frozen-plan LOS scenarios do not measure that end-to-end benefit.

Implementation verification checks case conservation, resource consistency,
information boundaries and reproducible scoring. It does not establish
clinical validity. Multi-day ward dynamics, emergency arrivals, patient
priority, unrestricted sequencing, clinically confirmed resource eligibility
and prospective workflow evaluation remain scientific extensions, rather
than demonstrated capabilities of the campaigns reported here.
